This paper identifies tangible design parameters that can lead to inaccurate estimates of relatively small effects. Low statistical power not only makes such effects difficult to detect but resulting significant estimates systematically exaggerate effect sizes on average. Through the literature on the short-term health effects of air pollution, we explore the prevalence of this issue and its policy implications. While some studies seem robust, others may suffer from substantial exaggeration. Real-data simulations highlight key drivers including the number of exogenous shocks, instrument strength, and outcome distribution. We propose approaches to evaluate and mitigate exaggeration when conducting non-experimental studies.